\section{Five Operations a Developer Cannot Name}
\label{sec:scenarios}

The five scenarios below come from our own use of \sgt{} while building it and from
watching colleagues work with coding agents on repositories we maintain. Each one
is an operation a developer tried to perform, stated in the words they used, and
all five fail for the same reason, that the developer names the work by the request
that produced it, because that is the level their theory of the codebase is pitched
at~\cite{naur1985}, while the record they have to operate through holds lines inside
commits (Figure~\ref{fig:granularity}). Every scenario has a workaround and all the
workarounds work, so the argument is about what the workarounds cost, and what they
cost is working out from the code which function served which request, a division the
developer stated in advance and nothing wrote down.

\subsection{``Commit this when it is done, and I do not know yet when that is''}

A developer working through an agent still has to choose when to commit, and in
this hour every moment available to them is wrong in a different way. Committing
after each request writes a version of \sym{api.py::handle} into the permanent
history that existed for twenty-nine minutes and that nobody ever ran the test
suite against, and it files the agent's rework of the cache under the retry commit,
because the agent did that rework while adding the retry. Committing once at the
end produces one commit holding all three requests and the rename, and this
developer split the difference by committing twice, which is the pair
Figure~\ref{fig:two-records} draws, each of them mixing two requests with the
rename. Staging hunk by hunk produces four clean
commits and asks the developer to decide, for each of fifteen hunks, which of those
four purposes it serves. Green and Petre named the first two costs premature
commitment, where a notation forces a decision before the information it needs
exists, and the third viscosity, the resistance a notation puts up once a decision
turns out to be wrong~\cite{green1996,blackwell2003}. Git asks for the
decomposition at the moment of writing, and in a session driven by requests the
boundary does exist at that moment: what the developer lacks is not the boundary,
which they stated themselves before any of the code existed, but the list of
functions on either side of it, which only the agent ever had.

\smallskip\noindent\emph{The operation.} Record the boundary the developer already
stated, without asking them to state it a second time in a different notation.

\subsection{``Change the rate limiter I built yesterday''}

The rate limiting from 14:02 used a fixed one-minute window, and three days later
it needs a sliding one. The work lands in two of the four functions the original
request produced, and the developer's first question is what the rate limiting
looked like before they touched it. Git can follow one function at a time:
\cmd{git log -L :handle:api.py} follows a function by matching a name pattern and
\cmd{git log -S} finds the commits where a string appeared or disappeared, so the
developer runs the first command four times, once per function they believe was
involved, and assembles the four answers themselves. The set has no name, so
nothing checks whether four was the right number. Codoban et al.\ asked developers
why they examine history and found the most common reason to be understanding the
rationale behind a piece of code~\cite{codoban2015}, and LaToza and Myers found
questions about why code is the way it is among the hardest developers report
answering~\cite{latoza2010}. The rationale in this case was one sentence of
English, typed three days ago into a window that is now closed.

\smallskip\noindent\emph{The operation.} Name the work one request produced, see
its current state and its previous state, and change it as a unit.

\subsection{``Take the caching back out and leave the rest alone''}

The caching turned out to serve a stale price for up to sixty seconds, and it has
to go. Its six edits are spread over three files and two commits, mixed with nine
hunks belonging to three other purposes, one of them adds a parameter to
\sym{db.py::query} while another passes it from \sym{api.py::handle}, and three
functions the rename rewrote afterwards call a helper the caching introduced.
\cmd{git revert} takes back a whole commit, so reverting either of these two
removes work the developer is keeping. The developer therefore checks out both
commits, selects the six hunks by hand, reverses them, and then confirms that
removing the parameter did not break \sym{handle} and that the three functions the
rename rewrote no longer call \sym{get\_cached}. The confirmation takes the time,
because the record contains no statement that the \sym{db.py::query} signature and
the \sym{handle} call site were changed for the same reason.

\smallskip\noindent\emph{The operation.} Remove a named set of edits, leave every
other edit to the same files in place, and name the code that still depends on the
removed set before anything is applied.

\subsection{``Send just the export work to the release branch''}

An export endpoint has to ship today, and it was written on a branch that also
carries a half finished retry policy. \cmd{git cherry-pick} moves a commit, and the
export work is two functions in a commit that also carries part of that retry
policy, so the usual answer is a temporary branch and an interactive rebase
that splits that commit, and the split has to be done again if more export work
arrives before the release. Darcs solves the part of this that is about lines,
since a Darcs patch carries the dependencies it was recorded with and moving it
involves no reapplication (Section~\ref{sec:related}). But the export code calls a
helper the rate limiter introduced, no line-level record states that call, and the
developer finds out when the release branch fails to import.

\smallskip\noindent\emph{The operation.} Carry a named set of edits and its
dependencies to another branch, and be told before the operation which dependencies
would come along.

\subsection{``Let someone review the rate limiter while I keep building on it''}

The rate limiter was finished at 14:30 and the caching that followed calls
\sym{api.py::\_bucket}, so the two form a stack of two layers with the rate limiter
underneath. Submitting them that way required the rate limiter to be on its own
branch at 14:30, which required knowing at 14:30 that the caching would build on
it. Stacked pull requests exist because review works best on small changes arriving
often~\cite{rigby2013} and because a reviewer's main difficulty is understanding
the change in front of them~\cite{bacchelli2013}, and Graphite and GitHub automate
every step after the decision, including rebasing the upper layer when a reviewer
asks for a change to the lower one (Section~\ref{sec:related}). The decision itself
belongs to the developer and has to be made before the work exists, and in an agent
session the layering becomes visible only once both layers are written. Green and
Petre's framework has a name for this, the hidden dependency, a relationship the
system relies on that the notation does not show, and Church et al.\ traced its
effects through version control in particular~\cite{green1996,church2014}. The
notation here is the branch, and a
branch says nothing about which functions call which.

\smallskip\noindent\emph{The operation.} Discover the layering after the work
exists, and submit one layer while continuing to work on the layer above it.

\subsection{What the five have in common}

Every one of these operations takes a set of edits as its argument, and the
developer identifies that set by the purpose the edits served. Git expresses a set
of lines, a set of files, and a set of commits, and a developer can build all five
operations out of those three. The price of building them is recovering a
decomposition that nobody recorded, and that recovery stayed affordable for as long
as the person paying it was the person who had written the code.
Section~\ref{sec:design} describes what we changed so that a developer can type the
argument these five operations want.
